\uuid{QTk3}
\titre{Arithmétique}
\auteur{bodin}
\organisation{exo7}
\licence{CC BY-NC-SA 4.0}
\datecreate{2020-01-21}
\niveau{L1}
\module{Algèbre}
\chapitre{Arithmétique dans Z}
\sousChapitre{Autre}
\difficulte{}
\theme{Arithmétique}
\exo7qcmid{arnaud.arnaud.013}

\contenu{
\begin{question}
\qclassification{103.99}
Quelles sont les assertions vraies ?
\begin{answers}
    \good{Le pgcd de $924$, $441$  et $504$   est $21$.}

    \bad{$627$ et $308$ sont premiers entre eux.}

    \bad{Si $p \ge 3$ est premier, alors $p!$ est premier.}

    \good{Soit $n \ge 2$ alors $n$ et $n+1$ sont premiers entre eux.}

    \good{Soit $n \ge 2$ un entier, le pgcd de $\{ in^i \text{ pour } i=1,\ldots ,100 \}$ est $n$.}
\end{answers}
\end{question}
}
